\documentclass[exercice]{thedocument}%
\usepackage{texmaths-tikz}%
\startdatakeys
\source{}
\cycle{}
\competencesDuSocle{}
\connaissancesRequises{}
\competencesTravaillees{}
\enddatakeys
\begin{document}%
    Nommer les angles ci-dessous de deux manières différentes.
    \vskip10pt\noindent
    \begin{tikzpicture}[scale=0.5]
        \AnglexOy[
            sommet name=\<S1;,
            right side name=\<r1;,
            left side name=\<l1;,
            color=Rainbow1-bleu,
            angle=105%
        ]
    \end{tikzpicture}\hskip30pt
    \begin{tikzpicture}[scale=0.5,rotate=80]
        \AnglexOy[
            sommet name=\<S2;,
            right side name=\<r2;,
            left side name=\<l2;,
            left side pos=below left,
            right side pos=right,
            color=Rainbow1-vert,
            angle=75%
        ]
    \end{tikzpicture}\vskip10pt
    \begin{tikzpicture}[scale=0.5,rotate=15]
        \AngleAOB[
            sommet name=\<S3;,
            right side name=\<r3;,
            left side name=\<l3;,
            left side pos=below left,
            color=Rainbow1-orange,
            angle=95%
        ]
    \end{tikzpicture}\hskip30pt
    \begin{tikzpicture}[scale=0.6,rotate=15]
        \AngleAOB[
            sommet name=\<S4;,
            right side name=\<r4;,
            left side name=\<l4;,
            left side pos=above left,
            right side pos=below right,
            color=Rainbow1-rouge,
            angle=25%
        ]
    \end{tikzpicture}
\end{document}
